% Generated draft: requires mathematical and source review.
\documentclass[11pt]{article}
\usepackage[margin=25mm]{geometry}
\usepackage{fontspec}
\setmainfont{Latin Modern Roman}
\usepackage{amsmath,amssymb,mathtools,unicode-math}
\setmathfont{Latin Modern Math}
\usepackage{enumitem,needspace,xurl,hyperref}
\setlength{\parindent}{0pt}
\setlength{\parskip}{5pt}
\setlength{\emergencystretch}{3em}
\newcommand{\N}{\mathbb{N}}
\newcommand{\Z}{\mathbb{Z}}
\newcommand{\Q}{\mathbb{Q}}
\newcommand{\R}{\mathbb{R}}
\newcommand{\C}{\mathbb{C}}
\newcommand{\F}{\mathbb{F}}
\newcommand{\A}{\mathbb{A}}
\newcommand{\PP}{\mathbb{P}}
\newcommand{\E}{\mathbb{E}}
\newcommand{\eps}{\varepsilon}
\newcommand{\dd}{\,\mathrm{d}}
\begin{document}
\section*{UnsolvedMath ID: 30003255}
\textbf{Source problem code:} catalogue entry
\subsection{Definitions and mathematical statement}
Let \((K,v)\) be a valued field, with valuation written additively as \(v:K\to\Gamma\cup\{\infty\}\), where \(\Gamma=v(K^\times)\) is an ordered abelian group and \(v(0)=\infty\). Write
\[
\mathcal O_v=\{a\in K:v(a)\geq 0\},\qquad
\mathfrak m_v=\{a\in K:v(a)>0\},\qquad
Kv=\mathcal O_v/\mathfrak m_v.
\]
Positive residue characteristic means \(\operatorname{char}(Kv)=p>0\); it does not require \(\operatorname{char}(K)>0\).

The valued field \((K,v)\) is \emph{extremal} if, for every integer \(n\geq 1\) and every polynomial \(f\in K[X_1,\ldots,X_n]\), the set
\[
\{v(f(a_1,\ldots,a_n)):(a_1,\ldots,a_n)\in\mathcal O_v^n\}
\subseteq\Gamma\cup\{\infty\}
\]
has a largest element, with \(\infty\) larger than every element of \(\Gamma\).

Use the standard convention that a \(\mathbb Z\)-group is an ordered abelian group elementarily equivalent to \((\mathbb Z,+,0,<)\) in the language of ordered abelian groups. A group \(\Gamma\) is \emph{divisible} if, for every \(\gamma\in\Gamma\) and every integer \(m\geq 1\), there exists \(\delta\in\Gamma\) such that \(m\delta=\gamma\). A field \(k\) is \emph{large} if every smooth geometrically integral curve over \(k\) having a \(k\)-rational point has infinitely many \(k\)-rational points. In characteristic \(p>0\), \(k\) is \emph{imperfect} if \(k^p=\{a^p:a\in k\}\ne k\).

An extension \((L,w)/(K,v)\) is \emph{immediate} if the induced inclusions of value groups and residue fields are equalities. The valued field \((K,v)\) is \emph{maximal} if it has no proper immediate valued field extension.

\textbf{Question concerning an additional assumption.} Consider valued fields \((K,v)\) satisfying \(\operatorname{char}(Kv)=p>0\) and
\[
\bigl[\Gamma\text{ is a }\mathbb Z\text{-group}\bigr]
\quad\text{or}\quad
\bigl[\Gamma\text{ is divisible and }Kv\text{ is large and imperfect}\bigr].
\]
Find a simple, useful additional assumption under which every such valued field is extremal. The source does not specify a class of admissible additional assumptions or mathematical criteria for being simple and useful; this part is therefore a request for a sufficient-condition criterion rather than a fully specified proposition.

\textbf{Question concerning maximal fields.} Is every maximal valued field \((K,v)\) satisfying
\[
\bigl[\Gamma\text{ is a }\mathbb Z\text{-group}\bigr]
\quad\text{or}\quad
\bigl[\Gamma\text{ is divisible and }Kv\text{ is large}\bigr]
\]
extremal? In this second question, the supplied wording omits both the imperfection requirement and an explicit repetition of positive residue characteristic. The imperfection requirement is therefore absent from the displayed question. It remains unresolved whether the quantifier in this second question ranges over all valued fields or only those with \(\operatorname{char}(Kv)>0\); these are distinct formulations requiring source clarification.
\subsection{Short English statement}
Ask for a simple, useful sufficient condition for extremality in positive residue characteristic when the value group is a \(\mathbb Z\)-group, or is divisible with large imperfect residue field. Also ask whether maximality suffices for a \(\mathbb Z\)-group, or for a divisible value group with large residue field. The characteristic scope of the second question is unclear. Extremality means that every polynomial attains a largest valuation on tuples from the valuation ring.
\subsection{Sources}
\begin{itemize}
\item Local UnsolvedMath registry, supplied complete source record, permanent ID \(30003255\), source problem code \texttt{catalogue entry}, in the Oberwolfach Reports Open Problems collection. The project identifies \texttt{data/unsolvedmath/problems.json} and \texttt{catalog/problems.json} as registry sources; the supplied record does not identify which file contains this entry.
\item Supplied report citation: \emph{Definability and Decidability Problems in Number Theory} (2017), Oberwolfach Reports, DOI 10.4171/owr/2016/49, pp. 2793--2866. \url{https://doi.org/10.4171/owr/2016/49}. The record associates the problem with 2016 and describes the report as documenting workshop open-problem discussions. No individual proposer or exact problem page is supplied.
\item Anscombe--Kuhlmann, \emph{Extremal valued fields}. \url{https://arxiv.org/abs/1407.3759}.
\item Kuhlmann, \emph{Extremal valued fields notes}. \url{https://fvkuhlmann.de/EXTRNOT.pdf}.
\item Literature assessment dated exactly \texttt{2026-08-21}: the supplied curation labels the problem open and states that there is no complete characterization of extremal valued fields in positive residue characteristic under the stated residue-field and value-group conditions. This is supplied dated curation derived from the OWR material, not an independently verified assessment. The supplied background expressly calls for independent verification of time-sensitive status and cited claims.
\item The registry labels the statement \texttt{corrected\_verified} and reports a manual semantic, typographic, and LaTeX review. These are supplied editorial assertions, not independent mathematical verification of this draft or of the cited sources.
\end{itemize}
\end{document}
